% !TeX root = ../../main.tex
\subsection{Protocol and Baselines}
Training and validation contain 72 and 36 scenes in 18 and 9 parameter groups.
Each group has four obstacle layouts: front blocked, rear blocked, both
blocked, and free. Nine templates vary door appearance, handle configuration,
friction, resistance, and geometry. Each final evaluation uses a new
18-group, 72-scene test with disjoint group and geometry identities.
The three test cohorts are named A, B, and C in collection order, with C the
latest cohort.
Validation determines checkpoints and decision thresholds before test
collection. The implementation and experiment records are available in
the project repository.\footnote{\url{https://github.com/blanknamezzz/mujoco-door-opening},
revision \texttt{0ae4c4d}.}

A scene is \emph{candidate-feasible} if at least one of its 16 physical
candidate records succeeds. Task correctness counts a successful opening in
a feasible scene or a refusal in a library-infeasible scene. Opening and
correct-refusal fractions use their respective scene counts as denominators.
We also measure collisions, false refusals, actual inspections, whole-episode
simulated duration, base path length, and net utility from
Equations~\ref{eq:utility}--\ref{eq:advantage}.

The primary comparison holds the frozen continuation and depth constraints
fixed while varying learned, always, and never inspection. For cohort C,
1152 physical candidate trials and 288 policy replays cover four
distinct deployed policies.
